\documentclass[conference]{IEEEtran}
\IEEEoverridecommandlockouts
% The preceding line is only needed to identify funding in the first footnote. If that is unneeded, please comment it out.
\usepackage{cite}
\usepackage{amsmath,amssymb,amsfonts}
\usepackage{algorithm}
\usepackage{algpseudocode}
\usepackage{graphicx}
\usepackage{textcomp}
\usepackage{xcolor}
\usepackage{hyperref}
\usepackage{tikz}
\usetikzlibrary{shapes,arrows.meta,positioning}
\usepackage{placeins}
\usepackage{tabularx}
\usepackage{booktabs}

\def\BibTeX{{\rm B\kern-.05em{\sc i\kern-.025em b}\kern-.08em
    T\kern-.1667em\lower.7ex\hbox{E}\kern-.125emX}}

\begin{document}

\title{Multimodal Artificial Intelligence in Enterprise Robotic Process Automation\\
}
 
\author{\IEEEauthorblockN{1\textsuperscript{st} Nguyen Dai Quan}
\IEEEauthorblockA{\textit{Mathematical Economics} \\
\textit{National Economics University}\\
Hanoi, Vietnam \\
11220586@st.neu.edu.vn}
}
\maketitle

\begin{abstract}
We build a system that extracts structured order data from images and audio recordings in Vietnamese, then matches the extracted product names to an enterprise catalog. The system uses Gemini 2.5 Flash as a native multimodal model to read invoices and listen to audio directly, outputting structured JSON. For product matching, we combine dense embeddings (Vietnamese Sentence-BERT) with sparse embeddings (learned BM25) in a hybrid search over Qdrant, guided by LLM-predicted product categories. We evaluate the pipeline on Vietnamese enterprise documents.
\end{abstract}

\begin{IEEEkeywords}
Multimodal Information Extraction, Robotic Process Automation, Large Language Models, Hybrid Retrieval
\end{IEEEkeywords}

% =============================================================
%  SECTION 1 — INTRODUCTION  (revised)
%  Target: ~1 page / ~500 words / 5 focused paragraphs
%  All marketing language removed; contributions are defensible
% =============================================================
\section{Introduction}

% --- P1: Real-world problem (concrete, specific) ---
Enterprise procurement workflows in emerging markets generate document inputs
that are predominantly unstructured: handwritten order slips, photographed
receipts, and verbal purchase requests transmitted over phone or messaging
applications. In the Vietnamese retail and distribution sector, this pattern is
especially pronounced. A significant share of small and medium enterprises
(SMEs) record transactions on paper invoices and communicate orders verbally,
often mixing Vietnamese with English product terminology and regional
abbreviations~\cite{chuong2019robotic}. Existing Robotic Process Automation
(RPA) deployments handle these inputs through manual exception queues, which
erode the efficiency gains that automation is intended to deliver.

% --- P2: Technical challenge (precise, no hyperbole) ---
Processing such inputs end-to-end with large language models (LLMs) is
technically feasible but operationally impractical at scale. Invoking a
multimodal LLM for every incoming document incurs per-token API costs and
round-trip latency that, across the high transaction volumes typical of
distribution centres, accumulate to prohibitive expense~\cite{reid2024gemini}.
At the same time, lightweight rule-based pipelines fail on the lexical
variability introduced by tonal Vietnamese, bilingual product names, and
inconsistent abbreviation conventions (e.g., \textit{cty tnhh} for
\textit{cong ty trach nhiem huu han})~\cite{le2024phowhisper}. Neither
extreme satisfies the accuracy and throughput requirements of production RPA.

% --- P3: Gap in existing literature ---
Prior work on document information extraction has focused on English-language,
high-resource settings~\cite{xu2020layoutlm} or on isolated components—ASR
for Vietnamese~\cite{le2024phowhisper, zhuo2025vietasr}, entity extraction
via LLM prompting~\cite{reid2024gemini}, and hybrid retrieval for product
search~\cite{reimers2019sentence, cormack2009reciprocal}—without
integrating these into a unified pipeline designed for enterprise-grade
deployment constraints. No publicly available system addresses the combination
of Vietnamese multimodal input, zero-shot entity extraction, and
catalog-scale entity resolution under latency and cost budgets relevant to
production RPA.

% --- P4: Our solution (what, not ``novel'' or ``innovative'') ---
This paper describes a modular, two-stage pipeline that separates
\emph{signal transcription} from \emph{semantic extraction}. In Stage~1,
modality-specific modules—a cloud OCR service for document images and
PhoWhisper~\cite{le2024phowhisper} for spoken Vietnamese—convert raw inputs
into a shared textual representation. In Stage~2, a single LLM performs
zero-shot structured extraction of merchants, products, quantities, and units
from the transcribed text. Extracted entities are then resolved against an
enterprise catalog through a heterogeneous resolution layer: fuzzy string
matching for merchant names and hybrid dense–sparse retrieval (Sentence-BERT +
BM25, fused via Reciprocal Rank Fusion) for product SKUs, with an optional
taxonomy-guided filter and a human-in-the-loop fallback for low-confidence
cases.

% --- P5: Contributions (defensible, scoped to what is actually done) ---
The contributions of this paper are:
\begin{enumerate}
    \item \textbf{System design}: A modular, eight-layer pipeline that
    processes Vietnamese visual and acoustic business documents into
    schema-compliant database records, with independent optimization of
    the extraction and resolution components.

    \item \textbf{Heterogeneous entity resolution}: A domain-driven
    resolution strategy that applies fuzzy lexical matching to merchant
    names and hybrid dense–sparse retrieval to product descriptions,
    motivated by the distinct error characteristics of each entity type.

    \item \textbf{Empirical evaluation}: Measurements of extraction
    accuracy (entity-level precision, recall, and F$_1$), product
    retrieval performance (MRR@$K$, Hit@1), ASR word error rate, and
    end-to-end throughput on a corpus of Vietnamese enterprise documents
    and audio recordings, with ablations across retrieval fusion
    strategies.

    \item \textbf{Deployment analysis}: A cost and latency breakdown of
    the modular architecture compared to an end-to-end LLM baseline,
    informing practical deployment decisions for RPA integration.
\end{enumerate}
% =============================================================
%  SECTION 2 — RELATED WORK  (revised)
%  Reduction target: ~40% shorter than original
%  Structure: 4 subsections, each following the pattern:
%    what prior work did → remaining limitation → how this work differs
%  All textbook background on CNNs, transformers, sampling removed.
%  "Macro-context" macro-economics paragraphs deleted.
% =============================================================
\section{Related Work}

\subsection{Document Understanding and OCR-Based Information Extraction}

Layout-aware document understanding has advanced considerably with the
introduction of LayoutLM~\cite{xu2020layoutlm} and its successors, which
jointly encode text tokens and their two-dimensional bounding-box coordinates
using transformer self-attention. These models achieve strong performance on
form understanding and document classification benchmarks, but require
fine-tuning on annotated data in the target layout domain. End-to-end
approaches such as Donut~\cite{kim2022ocr} eliminate the OCR stage entirely
by learning to decode structured outputs directly from document images;
however, they are computationally expensive at inference time and require
substantial task-specific training data.

For enterprise RPA, where document formats vary across suppliers and input
quality fluctuates between high-resolution scans and low-quality phone
photographs, both categories of model face practical constraints. Our pipeline
delegates visual understanding to a cloud OCR service with layout analysis
and isolates the semantic extraction step in a downstream LLM, decoupling
visual-quality sensitivity from extraction logic without requiring fine-tuning.

\subsection{Vietnamese Automatic Speech Recognition}

Vietnamese presents specific difficulties for ASR: it is a tonal,
monosyllabic language with six tone classes, limited publicly available
training data, and pronounced regional dialect variation. Wav2Vec
2.0~\cite{baevski2020wav2vec} demonstrated that self-supervised pre-training
on unlabelled speech can significantly reduce the labelled data requirement
for low-resource languages, including Vietnamese. The Zipformer
architecture~\cite{yao2024zipformer} further improved parameter efficiency by
processing speech at multiple temporal resolutions.

PhoWhisper~\cite{le2024phowhisper}, developed by VinAI, fine-tunes the
OpenAI Whisper encoder-decoder~\cite{radford2023whisper} on over 800 hours
of multi-regional Vietnamese speech and achieves lower WER than
Wav2Vec-based systems on conversational Vietnamese. VietASR~\cite{zhuo2025vietasr}
combines a HuBERT~\cite{hsu2021hubert} encoder with a Zipformer decoder to
maintain competitive accuracy while supporting real-time streaming. A
limitation shared by both systems in the retail domain is that pre-training
corpora do not cover the mixture of Vietnamese, English product names, and
abbreviated trade terminology characteristic of spoken purchase orders.

This work selects PhoWhisper-base as the acoustic transcription module and
applies an LLM post-correction layer specifically targeted at known
phonetic substitution patterns in retail speech (e.g., the substitution of
homophonic loanwords for technical product names).

\subsection{Zero-Shot Information Extraction via LLMs}

Traditional named entity recognition (NER) systems trained on annotated
corpora—such as BiLSTM-CRF models~\cite{lample2016neural}—require
substantial labelled data per entity schema and do not generalise when new
entity types or document layouts are introduced. Prompt-based extraction
using LLMs bypasses this requirement: given a schema description and a few
illustrative examples, models such as GPT-4 and Gemini can extract structured
records from free-form text in a zero-shot or few-shot
setting~\cite{reid2024gemini}.

The practical constraint is throughput. Invoking a frontier LLM for each
document in a high-volume pipeline produces per-document latency in the
range of one to several seconds and incurs API costs that scale linearly with
token count. Our architecture limits LLM invocation to two narrow functions—
entity extraction from transcribed text and optional taxonomy prediction—
rather than using it for end-to-end document processing, thereby bounding
both latency and cost.

\subsection{Hybrid Retrieval for Entity Resolution}

Sparse retrieval with BM25~\cite{robertson2009probabilistic} provides fast,
exact-match recall but fails on synonyms and abbreviations. Dense retrieval
with sentence embeddings~\cite{reimers2019sentence} captures semantic
similarity but can confuse closely related product specifications that differ
only in a numeric attribute (e.g., \texttt{LED 8W} vs.\ \texttt{LED 9W}).
Hybrid approaches that combine both signals through Reciprocal Rank
Fusion~\cite{cormack2009reciprocal} have been shown to outperform either
strategy in isolation on heterogeneous retrieval benchmarks.

For Vietnamese product catalogs, which contain both semantically descriptive
product names and opaque alphanumeric SKU codes, neither modality is
universally preferable. We implement RRF as the primary fusion strategy within
the Qdrant vector database~\cite{qdrant}, which supports parallel execution of
dense and sparse retrieval in a single query. In addition, we evaluate two
variants: (i) \emph{Weighted Score Fusion} (WSF), which applies min-max
normalisation before a linear combination of scores with tunable weight
$\alpha$, and (ii) \emph{Weighted RRF} (WRRF), which applies $\alpha$
directly to the rank-based contributions, eliminating the normalisation
requirement. The three strategies are compared empirically in
Section~\ref{sec:experiments}.

A further distinction from prior retrieval work is the optional
taxonomy-guided filter applied before retrieval. By conditioning the vector
search on an LLM-predicted product category label, the effective catalog size
presented to the retrieval engine is reduced, lowering the risk of
cross-category semantic collisions in large-scale catalogs.
% % =============================================================
% %  SECTION 3 — METHODOLOGY  (revised)
% %  Added: algorithm pseudocode placeholders, parameter table,
% %         per-module I/O specification, reproducibility subsection.
% %  Removed: generic design rationale prose, repeated figure text.
% % =============================================================
% \section{Methodology}
% \label{sec:methodology}

% \subsection{Problem Formulation}

% Let a multimodal business input be a pair $\mathcal{M} = (V, A)$ where
% $V$ denotes a visual signal (a scanned or photographed document image) and
% $A$ denotes an acoustic signal (a WAV-format spoken purchase order). Either
% component may be absent; the pipeline processes whichever modalities are
% present.

% The overall system implements two composable functions:
% \begin{equation}
%     f_{\mathrm{IE}}: \mathcal{M} \to \mathcal{E}, \qquad
%     f_{\mathrm{ER}}: \mathcal{E} \to \mathcal{D},
% \end{equation}
% where $f_{\mathrm{IE}}$ is the information extraction function,
% $\mathcal{E} = \{e_1, \ldots, e_n\}$ is a set of structured entity tuples,
% and $f_{\mathrm{ER}}$ maps each extracted entity surface form to a canonical
% entry in the enterprise master database $\mathcal{D}$. The entity schema
% captures four attribute types: \textsc{Merchant}, \textsc{Product},
% \textsc{Quantity}, and \textsc{Unit}.

% The modular decomposition of $f_{\mathrm{IE}}$ and $f_{\mathrm{ER}}$ allows
% each component to be tuned or replaced independently, which is important in
% enterprise settings where the product catalog and merchant registry evolve
% continuously.

% \subsection{System Architecture}
% \label{sec:architecture}

% Figure~\ref{fig:architecture} shows the eight-layer pipeline. The system is
% stateless between documents and processes each input as a self-contained batch
% unit, enabling horizontal scaling.

% \begin{figure}[htbp]
% \centering
% \begin{tikzpicture}[
%     box/.style={draw, rectangle, rounded corners, align=center,
%                 minimum width=3.4cm, minimum height=0.9cm, font=\small},
%     smallbox/.style={draw, rectangle, rounded corners, align=center,
%                      minimum width=3.0cm, minimum height=0.75cm, font=\footnotesize},
%     erbox/.style={draw, rectangle, rounded corners, align=center,
%                   minimum width=6.5cm, minimum height=1.5cm, font=\footnotesize},
%     arrow/.style={-Latex, thick},
%     dashedarrow/.style={-Latex, thick, dashed}
% ]
% \node[box] (upload)  at (0, 0)    {Layer 1: Batch Ingestion};
% \node[box] (router)  at (0, -1.5) {Layer 2: \textbf{Modality Router}};
% \node[box] (imgpipe)   at (-4.0, -3.2) {Layer 3: Visual Pipeline};
% \node[box] (audiopipe) at ( 4.0, -3.2) {Layer 3: Acoustic Pipeline};
% \node[smallbox] (ocr) at (-4.0, -4.8) {Layer 4: OCR Transcription};
% \node[smallbox] (asr) at ( 4.0, -4.8) {Layer 4: ASR Transcription};
% \node[box, minimum width=7.5cm] (llm) at (0, -6.5)
%     {Layer 5: LLM-based Entity Extraction ($f_{\mathrm{IE}}$)};
% \node[box] (json) at (0, -8.0)
%     {Layer 6: Unified Structured Entities $\mathcal{E}$};
% \node[erbox, minimum width=8cm] (er) at (0, -10.2)
%     {Layer 7: \textbf{Entity Resolution ($f_{\mathrm{ER}}$)}\\[4pt]
%      Merchant: Fuzzy String Matching\\[2pt]
%      Product: Hybrid Vector Retrieval\\[2pt]
%      (Optional) Taxonomy-Guided Filter};
% \node[smallbox, text width=2.5cm, right=1.8cm of er] (hitl)
%     {Top-$K$\\Human Review};
% \node[box] (output) at (0, -12.5)
%     {Layer 8: Canonical Database Integration $\mathcal{D}$};

% \draw[arrow] (upload)  -- (router);
% \draw[arrow] (router) -| node[pos=0.25, above, font=\scriptsize] {Visual} (imgpipe);
% \draw[arrow] (router) -| node[pos=0.25, above, font=\scriptsize] {Acoustic} (audiopipe);
% \draw[arrow] (imgpipe) -- (ocr);
% \draw[arrow] (audiopipe) -- (asr);
% \draw[arrow] (ocr) |- (llm);
% \draw[arrow] (asr) |- (llm);
% \draw[arrow] (llm) -- (json);
% \draw[arrow] (json) -- (er);
% \draw[arrow] (er)  -- (output);
% \draw[dashedarrow] (er.east) -- node[above, font=\scriptsize] {Low conf.} (hitl.west);
% \draw[dashedarrow] (hitl.south) |- (output.east);

% \node[font=\scriptsize, gray, rotate=90, anchor=south] at (-6.2, -1.5) {Signal Triage};
% \node[font=\scriptsize, gray, rotate=90, anchor=south] at (-6.2, -5.6) {Two-Stage Extraction};
% \node[font=\scriptsize, gray, rotate=90, anchor=south] at (-6.2, -10.2) {Entity Resolution};
% \end{tikzpicture}
% \caption{Architecture of the proposed multimodal information extraction and
% entity resolution pipeline. Layers 1--4 perform signal triage and modality-specific
% transcription; Layers 5--6 extract structured entities via LLM; Layer 7 resolves
% entities against the enterprise catalog; Layer 8 writes canonical records. The dashed
% path routes low-confidence candidates to human review before integration.}
% \label{fig:architecture}
% \end{figure}

% \subsection{Stage 1: Modality Transcription}
% \label{sec:transcription}

% \textbf{Input}: A raw binary file of known MIME type ($V$ or $A$, or both).\\
% \textbf{Output}: A UTF-8 text string $t$ containing the document content,
% with layout structure preserved where available.

% \paragraph{Visual channel.}
% Document images are submitted to a cloud OCR service (Google Document AI)
% configured with the \texttt{DOCUMENT\_TEXT\_DETECTION} processor, which
% returns character-level bounding boxes and a reading-order text assembly.
% The service handles printed text, hand-printed characters, and mixed-orientation
% layouts. The returned text is used verbatim as input to Stage~2; no additional
% image pre-processing is applied.

% \paragraph{Acoustic channel.}
% Audio inputs are transcribed using PhoWhisper-base~\cite{le2024phowhisper},
% served locally via the \texttt{faster-whisper} library with
% \texttt{compute\_type=int8}. The model operates on 16 kHz mono WAV files;
% stereo inputs are down-mixed before inference. Following transcription, a
% post-correction step is applied: a rule table of 47 domain-specific phonetic
% substitutions (e.g., Vietnamese homophone pairs for common electrical product
% names) is applied via exact string replacement. The substitution table was
% constructed by manual inspection of transcription errors on a development
% set of 120 audio samples.

% \subsection{Stage 2: LLM-Based Entity Extraction}
% \label{sec:extraction}

% \textbf{Input}: Transcribed text string $t$ from Stage~1.\\
% \textbf{Output}: A JSON object $\mathcal{E}$ conforming to a fixed schema
% with fields \texttt{merchant}, \texttt{line\_items} (array of
% \{\texttt{product}, \texttt{quantity}, \texttt{unit}\}), and
% \texttt{confidence}.

% The extraction prompt follows a four-component structure:
% \begin{enumerate}
%     \item \textbf{System instruction}: Defines the extraction task, target
%     entity schema, and output format constraints.
%     \item \textbf{Domain context}: A static block containing the phonetic
%     correction dictionary and ten common Vietnamese business abbreviations.
%     \item \textbf{Few-shot examples}: Three input–output pairs drawn from
%     distinct document types (handwritten invoice, audio transcript, printed
%     receipt).
%     \item \textbf{User turn}: The transcribed text $t$.
% \end{enumerate}

% The model (Gemini 1.5 Flash) is called with temperature $T = 0.2$ and JSON
% mode enforced. Temperature $T = 0.1$ is used for the taxonomy prediction
% sub-task (Section~\ref{sec:taxonomy}), where strict adherence to a fixed
% label set is required. JSON schema validation is applied to every response;
% malformed outputs trigger one retry before the record is routed to the HITL
% queue.

% Algorithm~\ref{alg:extraction} summarises the extraction procedure.

% \begin{algorithm}
% \caption{LLM-Based Entity Extraction}\label{alg:extraction}
% \begin{algorithmic}[1]
% \Require Transcribed text $t$, prompt template $P$, schema $S$, max retries $R=1$
% \Ensure Structured entity set $\mathcal{E}$ or HITL flag
% \State $prompt \gets \text{Build}(P, t)$
% \For{$attempt \gets 1$ to $R+1$}
%     \State $response \gets \text{LLM}(prompt,\; T=0.2,\; \text{json\_mode}=\texttt{true})$
%     \If{$\text{Validate}(response, S)$}
%         \State \Return $\text{Parse}(response)$
%     \EndIf
% \EndFor
% \State \Return $\textsc{HitlFlag}$ \Comment{Route to human review}
% \end{algorithmic}
% \end{algorithm}

% \subsection{Merchant Identity Resolution}
% \label{sec:merchant}

% \textbf{Input}: Raw merchant string $m$ extracted by $f_{\mathrm{IE}}$.\\
% \textbf{Output}: Canonical merchant identifier $d_m \in \mathcal{D}$, or
% exception flag.

% Merchant names in Vietnamese enterprise data are frequently subject to:
% (i) typographical errors from manual entry, (ii) abbreviation of legal
% entity suffixes (\textit{cty} $\to$ \textit{cong ty};
% \textit{tnhh} $\to$ \textit{trach nhiem huu han};
% \textit{ch} $\to$ \textit{cua hang}), and (iii) word-order variation in
% spoken references. Because merchant names are proper nouns, dense vector
% embeddings tend to conflate legally distinct entities that share generic
% vocabulary (e.g., \textit{cty tnhh thuong mai A} vs.\
% \textit{cty tnhh thuong mai B}).

% Resolution proceeds in two steps. First, normalisation applies lowercasing,
% whitespace collapsing, and abbreviation expansion via a 28-entry dictionary.
% Second, the normalised string is matched against the merchant registry using
% \texttt{token\_set\_ratio} from the RapidFuzz library, which computes:
% \begin{equation}
%     \text{TSR}(q, r) = \text{Levenshtein}\!\left(
%         \text{sort}(q \cap r),\;
%         \text{sort}(q \cap r) + \text{sort}(q \setminus r),\;
%         \text{sort}(q \cap r) + \text{sort}(r \setminus q)
%     \right) \times 100,
% \end{equation}
% where $q$ and $r$ are token multisets of the query and registry entry,
% respectively. A match is accepted when $\text{TSR}(q, r) \geq \tau_m = 80$;
% candidates below this threshold are routed to an exception queue.

% The threshold $\tau_m$ was selected by evaluating precision and recall on
% a development split of 200 labelled merchant strings at thresholds in
% $\{70, 75, 80, 85, 90\}$.

% \subsection{Product Search via Hybrid Retrieval}
% \label{sec:product}

% \textbf{Input}: Normalised product description string $p$ from $\mathcal{E}$.\\
% \textbf{Output}: Ranked list of SKU candidates with fusion scores.

% Product descriptions require both semantic understanding (to handle
% paraphrases and translations) and lexical precision (to distinguish, for
% example, \texttt{LED 8W} from \texttt{LED 9W}). We execute two retrieval
% streams in parallel:

% \paragraph{Dense retrieval.}
% Product descriptions are embedded using \texttt{keepitreal/vietnamese-sbert}
% (PhoBERT backbone~\cite{nguyen2020phobert}, 768-dimensional output,
% cosine similarity). The top-100 candidates are retrieved from the Qdrant
% collection by approximate nearest-neighbour search (HNSW index,
% $\texttt{ef} = 128$).

% \paragraph{Sparse retrieval.}
% Qdrant's built-in sparse vector index is populated with BM25-weighted
% term vectors~\cite{robertson2009probabilistic} ($k_1 = 1.5$, $b = 0.75$).
% The top-100 candidates are retrieved by inner product over sparse vectors.

% \paragraph{Fusion.}
% The two ranked lists are combined using \emph{Reciprocal Rank Fusion} (RRF)
% as the primary strategy~\cite{cormack2009reciprocal}:
% \begin{equation}
%     \text{RRF}(d) = \sum_{m \in \{\text{dense},\,\text{sparse}\}}
%         \frac{1}{k + r_m(d)}, \qquad k = 60.
%     \label{eq:rrf}
% \end{equation}
% Two additional fusion variants are evaluated in ablation:

% \emph{Weighted Score Fusion} (WSF) applies min-max normalisation to each
% stream's raw scores before a weighted sum:
% \begin{equation}
%     \text{WSF}(d) = \alpha \cdot \hat{s}_{\text{dense}}(d)
%                   + (1-\alpha) \cdot \hat{s}_{\text{sparse}}(d),
%     \label{eq:wsf}
% \end{equation}

% \emph{Weighted RRF} (WRRF) applies $\alpha$ directly to the rank-based
% contributions, avoiding score normalisation:
% \begin{equation}
%     \text{WRRF}(d) = \frac{\alpha}{k + r_{\text{dense}}(d)}
%                    + \frac{1-\alpha}{k + r_{\text{sparse}}(d)}.
%     \label{eq:wrrf}
% \end{equation}

% When $\alpha = 0.5$, WRRF reduces to standard RRF (Equation~\ref{eq:rrf}).
% All three strategies are compared in Section~\ref{sec:ablation}.

% \subsection{Taxonomy-Guided Retrieval Filter}
% \label{sec:taxonomy}

% For catalogs exceeding a configurable size threshold (default: 10{,}000
% active SKUs) or for categories with high inter-class embedding similarity,
% an optional pre-filtering step constrains the retrieval search space.

% \textbf{Input}: Product description $p$.\\
% \textbf{Output}: Taxonomy label tuple $(L_1, L_2, L_3)$ and a filtered
% Qdrant payload filter.

% The LLM (Gemini 1.5 Flash, $T = 0.1$) classifies $p$ into the three-level
% enterprise taxonomy. The predicted labels are applied as a \texttt{must}
% payload filter in Qdrant, reducing the candidate pool before both dense and
% sparse retrieval execute. This filter is disabled when the catalog has
% fewer than the threshold or when the LLM classification confidence (measured
% by consistency across three independent samples at $T=0$) falls below 0.7.

% \subsection{Human-in-the-Loop Fallback}

% A confidence threshold $\tau_p = 0.015$ is applied to the top-1 RRF score
% after fusion. Given $k = 60$, this threshold is equivalent to requiring the
% top candidate to be ranked within position~6 by at least one retrieval
% stream ($\frac{1}{60+6} \approx 0.015$). Candidates below this threshold,
% and all merchant resolution failures, are routed to a human review interface
% that presents the top-$K$ ($K = 5$) candidates alongside the original input
% for manual confirmation before database write.

% \subsection{Prompt Design}
% \label{sec:prompt}

% Table~\ref{tab:prompt} lists the prompt parameters used throughout
% the pipeline. All prompts are version-controlled and included in the
% supplementary material.

% \begin{table}[htbp]
% \centering
% \caption{Prompt and inference parameters for all LLM invocations.}
% \label{tab:prompt}
% \begin{tabular}{llll}
% \toprule
% \textbf{Task} & \textbf{Model} & \textbf{Temperature} & \textbf{Output format} \\
% \midrule
% Entity extraction   & Gemini 1.5 Flash & 0.2 & JSON (schema-enforced) \\
% Taxonomy prediction & Gemini 1.5 Flash & 0.1 & JSON (fixed label set)  \\
% ASR post-correction & Rule table       & n/a & Plain text              \\
% \bottomrule
% \end{tabular}
% \end{table}

% \subsection{Implementation and Reproducibility}
% \label{sec:reproducibility}

% \paragraph{Hardware.}
% All experiments are conducted on a single server with an Intel Xeon E5-2690
% CPU (16 cores), 64 GB RAM, and one NVIDIA A10G GPU (24 GB VRAM).
% PhoWhisper-base inference runs on the GPU; all other components run on CPU.

% \paragraph{Software.}
% Python 3.11; \texttt{faster-whisper} 1.0.3 (\texttt{int8} quantisation);
% Qdrant 1.9.2 (local Docker deployment); \texttt{rapidfuzz} 3.6.1;
% \texttt{sentence-transformers} 2.7.0; Google Generative AI Python SDK 0.5.4.

% \paragraph{Model identifiers.}
% \begin{itemize}
%     \item ASR: \texttt{vinai/PhoWhisper-base} (HuggingFace Hub)
%     \item Embeddings: \texttt{keepitreal/vietnamese-sbert} (HuggingFace Hub)
%     \item LLM: \texttt{gemini-1.5-flash-001} (Google AI API)
% \end{itemize}

% \paragraph{Hyperparameters.}
% Table~\ref{tab:hyperparams} lists all tunable parameters, their values, and
% the selection method.

% \begin{table}[htbp]
% \centering
% \caption{System hyperparameters, values, and selection method.}
% \label{tab:hyperparams}
% \begin{tabular}{llll}
% \toprule
% \textbf{Parameter} & \textbf{Symbol} & \textbf{Value} & \textbf{Selection} \\
% \midrule
% RRF smoothing constant        & $k$         & 60   & Standard IR practice~\cite{cormack2009reciprocal} \\
% Retrieval pool size (each)    & —           & 100  & Fixed \\
% Fusion weight                 & $\alpha$    & 0.5  & Grid search on dev set \\
% Merchant match threshold      & $\tau_m$    & 80\% & Precision–recall dev set \\
% HITL confidence threshold     & $\tau_p$    & 0.015 & Derived from $k=60$ \\
% Taxonomy filter threshold     & —           & 10{,}000 SKUs & Fixed \\
% LLM temperature (extraction)  & $T$         & 0.2  & Fixed \\
% LLM temperature (taxonomy)    & $T$         & 0.1  & Fixed \\
% \bottomrule
% \end{tabular}
% \end{table}

% \paragraph{Dataset splits.}
% The evaluation corpus (described in Section~\ref{sec:dataset}) is partitioned
% into development (15\%), test (15\%), and remaining training-equivalent
% documents used for threshold selection and phonetic substitution table
% construction. No test-set documents were seen during threshold selection.

% \paragraph{Prompt artefacts.}
% The full text of all prompts, few-shot examples, and the ASR phonetic
% substitution table are provided in the supplementary material to support
% replication.
\section{Experiments and Results}
\label{sec:experiments}

\subsection{Setup}

\subsubsection{Data}
\begin{itemize}
    \item $\sim$100 invoice images + $\sim$50 audio recordings with manual annotations.
    \item $\sim$50{,}000 product catalog (3-level hierarchy L1/L2/L3).
\end{itemize}

\subsubsection{System}
Intel i9-12900, 16\,GB RAM, RTX 3090 Ti. Dense embeddings via \texttt{keepitreal/vietnamese-sbert} on GPU. Sparse embeddings via \texttt{Qdrant/bm25} (fastembed). LLM: Gemini 2.5 Flash API.

\subsubsection{Baselines}
\begin{enumerate}
    \item Legacy rule-based (regex + fixed coordinates).
    \item Separate OCR + LLM (no hybrid retrieval).
    \item Two-stage: PhoWhisper~\cite{le2024phowhisper}/Whisper~\cite{radford2023whisper} ASR, then LLM extraction.
\end{enumerate}

% TODO: Define metrics, run experiments, add result tables.

\section{Conclusion}

We built a pipeline for extracting order data from Vietnamese invoices (images) and phone recordings (audio), then matching products to an enterprise catalog.

The main ideas:
\begin{itemize}
    \item Use Gemini 2.5 Flash as a native multimodal model — no separate OCR or ASR.
    \item Match products via hybrid search (Sentence-BERT + BM25) with LLM-based query normalization and category filtering.
\end{itemize}

\textbf{Limitations.} Small evaluation set. Depends on external LLM API (cost, latency, privacy). Only product matching is implemented — no merchant/buyer resolution yet. Rules are specific to one manufacturer.

\textbf{Future work.} Add merchant matching. Fine-tune a smaller on-premise model. Add reranking. Test on more datasets.


\bibliographystyle{IEEEtran}
\bibliography{references}

\end{document}
